\section{Case Study 2: Kimi K1.5}

Kimi K1.5 和 R1 几乎是同时代的另一个重要例子。课堂把它挑出来，是因为它和 R1 有相似的目标，但在细节上提供了不同的工程视角。

\begin{figure}[H]
\centering
\includegraphics[width=0.95\textwidth]{slides-images/slide-051.jpg}
\caption{Kimi K1.5：另一个用 RL 把 reasoning 推上去的开放模型}
\end{figure}

\subsection{数据、长度控制和系统细节}

Kimi 这条线的重点很实际：
\begin{itemize}
    \item 先做难度过滤和数据课程安排。
    \item 先用长 CoT 的 SFT 打底。
    \item 再用自己的 policy gradient / RL 配方去做强化。
    \item 训练中还加入了长度控制，避免 CoT 无限制拉长。
\end{itemize}

课堂特别强调，RL 系统本身很吃工程：on-policy rollouts 慢、训练框架切换麻烦、长 CoT 让 batch 更不均匀。这些都不是“算法论文里一个公式”能解决的。

\begin{figure}[H]
\centering
\includegraphics[width=0.95\textwidth]{slides-images/slide-052.jpg}
\caption{Kimi 的长 CoT reasoning strategy：数据、SFT、RL 三步走}
\end{figure}

\begin{figure}[H]
\centering
\includegraphics[width=0.95\textwidth]{slides-images/slide-053.jpg}
\caption{Kimi 的数据整理和 SFT：去掉太容易、太噪声、或格式不合适的样本}
\end{figure}

\begin{figure}[H]
\centering
\includegraphics[width=0.95\textwidth]{slides-images/slide-054.jpg}
\caption{Kimi 的 RL：仍然是某种基于偏好或可验证 reward 的策略优化}
\end{figure}

\begin{figure}[H]
\centering
\includegraphics[width=0.95\textwidth]{slides-images/slide-055.jpg}
\caption{长度控制：让更长的 rollout 不是无条件地更受鼓励}
\end{figure}

\begin{figure}[H]
\centering
\includegraphics[width=0.95\textwidth]{slides-images/slide-057.jpg}
\caption{RL infra 的重要性：rollout、框架切换、长序列不均匀性都会影响训练效率}
\end{figure}

